\chapter{Discussion}
\label{chap:discussion}

This chapter interprets the results of \Cref{chap:results} in the broader
context of the \HZIED framework and the \PLATO mission.


\section{The Dominant Failure Mode}
\label{sec:failure_mode}

\emph{Scaffold.}
Which physical effect most strongly erodes the \HZIED signal under
\PLATO-realistic assumptions: composition diversity, finite steam-atmosphere
lifetime, or the interplay of the two. Interpretation in terms of the
underlying planetary evolution.


\section{Implications for Mission Design}
\label{sec:mission_design}

\emph{Scaffold.}
Sample size versus per-planet precision trade-off in view of the
detectability boundary. Requirements on ground-based radial-velocity
follow-up. Opportunity for the \PLATO Guest Observer programme to extend the
hot-side sample.


\section{Limitations}
\label{sec:limitations}

\emph{Scaffold.}
Parametric atmospheric lifetime rather than integrated XUV escape modelling.
Fixed Bond albedo across the population. Absence of three-dimensional
atmospheric dynamics. Single $S_{\mathrm{thresh}}$ independent of host-star
spectral type. Assumed rocky mass--radius baseline.


\section{Alternative Interpretations}
\label{sec:alternatives}

\emph{Scaffold.}
If a step in mean radius is detected in real \PLATO data, what alternatives
to the \HZIED must be excluded before the interpretation is secure? The
cosmic shoreline of \textcite{OwenWu2017}, migration-driven compositional
gradients, and residual H/He envelopes on close-in sub-Neptunes.
